\subsection{Strong Typing in Operations}

Python names are dynamically bound, but Python operations are strongly checked at runtime. Dynamic typing does not mean ``anything combines with anything.''

\subsubsection{Why Heterogeneous Operations Such as \texttt{"3" + 4} Fail Instead of Silently Converting}

The expression \texttt{"3" + 4} raises \texttt{TypeError}. For integers, \texttt{+} means numeric addition via \texttt{\_\_add\_\_}. For strings, \texttt{+} means concatenation. The mixed expression is ambiguous: should Python produce \texttt{7}, \texttt{"34"}, or convert one operand silently?

Python refuses to guess. This is strong typing in practice:

\begin{quote}
Python does not silently convert unrelated object types to make an operation succeed.
\end{quote}

Both \texttt{str} and \texttt{int} implement \texttt{\_\_add\_\_}, but not for each other as partners. String repetition \texttt{"D" * 3} works because \texttt{str.\_\_mul\_\_} accepts \texttt{int}. String subtraction \texttt{"D" - 1} fails because no such protocol exists.

Strong typing does not forbid all mixed-type operations. It requires that mixed-type operations be explicitly defined by the involved types---as in the numeric tower, not as arbitrary cross-domain coercion.

\subsubsection{Explicit Conversion with \texttt{int()}, \texttt{float()}, \texttt{complex()}, \texttt{bool()}, and \texttt{str()}}

Crossing type boundaries requires explicit conversion:

\begin{verbatim}
int("3") + 4          # 7
"3" + str(4)          # "34"
\end{verbatim}

These produce different results because they express different intentions. Built-in converters obey target-type rules: \texttt{int("3.14")} raises \texttt{ValueError}; \texttt{int(float("3.14"))} truncates toward zero. \texttt{bool("False")} is \texttt{True} because non-empty strings are truthy---\texttt{bool()} tests truthiness, it does not parse English.

\subsubsection{Numeric Promotion Boundaries: Integer, Float, Complex, and Boolean Interactions}

Python defines mixed operations within its numeric tower:

\begin{verbatim}
int + int       -> int
int + float     -> float
bool + int      -> int      (True counts as 1)
bool + float    -> float
\end{verbatim}

These are defined interactions, not silent conversion of unrelated domains. \texttt{42 + "1"} raises \texttt{TypeError}. \texttt{True + 41 == 42} succeeds because \texttt{bool} is an \texttt{int} subclass---a legacy design choice, not arbitrary coercion.

Equality across numeric types is defined (\texttt{3.0 == 3}, \texttt{True == 1}); ordering of complex numbers against reals is not (\texttt{(3+4j) < 10} raises \texttt{TypeError}).

The operational summary:

\begin{verbatim}
"3" + 4          -> TypeError
int("3") + 4     -> 7
42 + 0.5         -> 42.5
True + 41        -> 42
\end{verbatim}

Names are flexible. Operations are checked. Objects keep their runtime type rules.
